%% kapa: exam=1 round=26 year=2015 period=1 form=A subject=economics subject_ko=경제원론 seq=34 number=74
%% source: Q-Net 감정평가사 기출문제 (cst00309) artlSeq=5138001 file=extracted/1008/5138001/제26회_감정평가사 1차_기출/제26회_감정평가사 1차_기출/1교시_감평_A형.pdf
%% conversion: kordoc 4.24.51 -> tools/gen1.py (auto); review=proofread-sonnet; flags=none
\begin{question}{74}
각 경제학파별 경제안정화정책에 관한 설명으로 옳지 않은 것은?
\begin{choices}
  \item 고전학파는 구축효과, 화폐의 중립성을 들어 경제안정화정책을 쓸 필요가 없다고 주장한다.
  \item 케인즈경제학자(Keynesian)는 IS곡선이 가파르고, LM곡선은 완만하므로 적극적인 재정정책이 경제안정화정책으로 바람직하다고 주장한다.
  \item 통화주의자(Monetarist)는 신화폐수량설, 자연실업률 가설을 들어 재량적인 경제안정화정책을 주장한다.
  \item 새고전학파(New Classical School)는 예상치 못한 경제안정화정책은 일시적으로 유효할 수 있다는 점을 인정한다.
  \item 새케인즈학파(New Keynesian School)는 임금과 물가가 경직적인 경우에는 경제안정화정책이 유효하다고 주장한다.
\end{choices}
\end{question}
\answer{3}
